\documentclass[../main.tex]{subfiles}

\begin{document}


\section{Khảo sát thực tế}

\subsection{Quy trình học lập trình hiện nay}

Theo báo cáo thị trường nhân lực công nghệ thông tin của TopDev, mỗi năm Việt Nam có hơn 57.000 sinh viên CNTT tốt nghiệp nhưng chỉ khoảng 30\% trong số đó đáp ứng được yêu cầu chuyên môn thực tế từ phía doanh nghiệp. Phần lớn sinh viên ra trường nắm lý thuyết tương đối vững nhưng hạn chế về khả năng thực hành và tư duy xử lý các tình huống phát sinh trong bối cảnh làm việc thực tế. Khoảng cách này bắt nguồn từ nhiều nguyên nhân, trong đó một yếu tố quan trọng là cách thức sinh viên tiếp cận lập trình trong giai đoạn học tập tại trường.

Quy trình học lập trình phổ biến hiện nay thường diễn ra theo một vòng lặp quen thuộc: tiếp nhận lý thuyết trên lớp hoặc qua video, tự ngồi làm bài tập, gặp lỗi, rồi tìm kiếm giải pháp qua Google hoặc các công cụ AI như ChatGPT. Vấn đề nằm ở chỗ mỗi lần tìm kiếm như vậy đều bắt đầu lại từ đầu, không có ngữ cảnh về những gì người học đã biết, đang gặp khó khăn ở đâu, hay đã mắc lỗi tương tự bao nhiêu lần trước đó. Các công cụ AI thông thường trả lời đúng câu hỏi được đặt ra nhưng hoàn toàn không lưu lại thông tin về người đặt câu hỏi. Điều này khiến sinh viên dù có thể tìm ra lời giải cho bài toán hiện tại nhưng lại chưa thực sự rèn luyện được khả năng tự phân tích và xử lý vấn đề. Vì vậy, các lỗi tương tự vẫn tiếp tục lặp lại ở những bài tập sau.

Ngoài ra, sĩ số lớp học lớn cũng tạo ra không ít khó khăn trong quá trình giảng dạy. Giảng viên khó theo dõi sát tiến độ của từng sinh viên, đồng thời khó điều chỉnh phương pháp hoặc kịp thời nhận ra khi người học bắt đầu mất kiến thức nền tảng. Sinh viên tiếp thu nhanh phải ngồi chờ, trong khi sinh viên tiếp thu chậm không dám hỏi vì sợ làm chậm cả lớp. Đây là bài toán cấu trúc không thể giải quyết chỉ bằng cách tăng số lượng giảng viên hay điều chỉnh giáo trình.

Từ thực trạng đó, xuất hiện nhu cầu xây dựng một hệ thống có thể hỗ trợ người học xuyên suốt quá trình học tập. Hệ thống không chỉ dừng lại ở việc trả lời câu hỏi mà còn cần theo dõi tiến độ học, phát hiện điểm còn hạn chế và đưa ra cách hỗ trợ phù hợp với từng sinh viên. Nội dung khảo sát ở phần tiếp theo sẽ phản ánh rõ vấn đề này thông qua ý kiến từ chính các sinh viên đang theo học lập trình.

\subsection{Kết quả khảo sát}

Nhằm tạo cơ sở thực tiễn cho quá trình phát triển hệ thống, em đã thực hiện một cuộc khảo sát trực tuyến vào tháng 3/2026. Kết quả thu được gồm 47 phản hồi hợp lệ từ sinh viên học lập trình đến từ 15 trường đại học khác nhau. Đối tượng tham gia chủ yếu là sinh viên chuyên ngành Công nghệ Thông tin, chiếm hơn 87\%, phần còn lại là sinh viên ngành khác có học lập trình như môn tự chọn hoặc tự học thêm ngoài chương trình.

\begin{figure}[H]
\centering
\includegraphics[width=1.0\textwidth]{khao_sat_chuyen_nganh_sinh_vien.png}
\caption{Khảo sát chuyên ngành sinh viên}
\label{fig:khao-sat-chuyen-nganh-sinh-vien}
\end{figure}

\subsubsection{Về kinh nghiệm và ngôn ngữ học}

Gần 51\% người tham gia có dưới sáu tháng kinh nghiệm lập trình, đây là nhóm mới bắt đầu, chưa hình thành thói quen tự học có hệ thống. Nhóm có trên hai năm kinh nghiệm chiếm 32\%, chủ yếu là sinh viên năm ba, năm tư đã qua các môn lập trình cơ sở. Ngôn ngữ phổ biến nhất là C++ với gần một nửa số phản hồi, tiếp theo là Python, Java và JavaScript, phản ánh đúng xu hướng đào tạo lập trình cơ sở tại các trường kỹ thuật Việt Nam hiện nay.

\begin{figure}[H]
\centering
\includegraphics[width=1.0\textwidth]{khao_sat_kinh_nghiem.png}
\caption{Khảo sát kinh nghiệm}
\label{fig:khao-sat-kinh-nghiem}
\end{figure}

\begin{figure}[H]
\centering
\includegraphics[width=1.0\textwidth]{khao_sat_ngon_ngu_hoc.png}
\caption{Khảo sát ngôn ngữ học}
\label{fig:khao-sat-ngon-ngu-hoc}
\end{figure}

\subsubsection{Về nguồn học và thói quen sau lý thuyết}

Khóa học trên lớp và YouTube là hai kênh được sử dụng nhiều nhất. Đáng chú ý là hỏi đáp AI đã trở thành nguồn học thứ ba, vượt qua cả tài liệu chính thức và nền tảng học trực tuyến trả phí. Điều này cho thấy sinh viên đã chủ động tích hợp AI vào quy trình học tập, song đây là AI tổng quát, không được tối ưu cho bối cảnh học lập trình cá nhân. Sau khi tiếp nhận lý thuyết, phần lớn người được hỏi chọn làm bài tập thực hành ngay, cho thấy nhận thức về tầm quan trọng của thực hành khá tốt. Vấn đề thực sự nằm ở giai đoạn sau khi gặp lỗi, không phải giai đoạn bắt đầu.

\begin{figure}[H]
\centering
\includegraphics[width=1.0\textwidth]{khao_sat_nguon_hoc.png}
\caption{Khảo sát nguồn học}
\label{fig:khao-sat-nguon-hoc}
\end{figure}

\begin{figure}[H]
\centering
\includegraphics[width=1.0\textwidth]{khao_sat_thoi_quen_sau_ly_thuyet.png}
\caption{Khảo sát thói quen sau lý thuyết}
\label{fig:khao-sat-thoi-quen-sau-ly-thuyet}
\end{figure}

\subsubsection{Về khó khăn trong quá trình học}

Đây là nhóm câu hỏi phản ánh rõ nhất khoảng trống mà hệ thống PPAT hướng đến lấp đầy. Gần 38\% người tham gia cho biết khó khăn lớn nhất là không biết bắt đầu từ đâu, tức là thiếu lộ trình học tập cụ thể và không có người dẫn dắt định hướng. Nhóm thứ hai chiếm 26\% gặp khó khăn khi không biết cách sửa lỗi code, đây là kỹ năng debug chỉ có thể hình thành qua thực hành lặp lại với phản hồi đúng thời điểm. Khoảng 15\% không hiểu được khái niệm lập trình dù đã đọc tài liệu, và 15\% thiếu người hướng dẫn bên cạnh khi cần. Tổng hợp lại, hơn 90\% người tham gia đang gặp ít nhất một trong những vấn đề mà một gia sư cá nhân có thể trực tiếp giải quyết.

\begin{figure}[H]
\centering
\includegraphics[width=1.0\textwidth]{khao_sat_kho_khan_trong_qua_trinh_hoc.png}
\caption{Khảo sát khó khăn trong quá trình học}
\label{fig:khao-sat-kho-khan-trong-qua-trinh-hoc}
\end{figure}

\subsubsection{Về hành vi khi gặp lỗi code}

Hơn một nửa số người được khảo sát, chiếm 68\%, chọn công cụ AI làm điểm dừng đầu tiên khi gặp lỗi, vượt xa so với tìm kiếm trên Google ở mức 13\% và hỏi bạn bè hoặc giảng viên ở mức 17\%. Con số này một mặt khẳng định AI đã trở thành công cụ debug mặc định trong cộng đồng sinh viên, mặt khác đặt ra câu hỏi về chất lượng hỗ trợ. AI tổng quát có thể xử lý lỗi cụ thể trong một phiên làm việc, nhưng không giải thích được tại sao sinh viên đó hay mắc lỗi đó, cũng không theo dõi xem lỗi có tái diễn ở những bài sau hay không.

\begin{figure}[H]
\centering
\includegraphics[width=1.0\textwidth]{khao_sat_hanh_vi_khi_gap_loi_code.png}
\caption{Khảo sát hành vi khi gặp lỗi code}
\label{fig:khao-sat-hanh-vi-khi-gap-loi-code}
\end{figure}

\subsubsection{Về mức độ sử dụng AI trong học tập}

68\% sử dụng AI thường xuyên và 26\% sử dụng thỉnh thoảng, tổng cộng hơn 94\% đã có thói quen tương tác với AI trong quá trình học. Đây là tín hiệu quan trọng cho thấy người dùng mục tiêu đã quen với giao diện chat-based và không có rào cản tâm lý khi sử dụng AI. Điều cần cải thiện không phải là thay đổi hành vi người dùng mà là nâng cấp chất lượng AI họ đang dùng, từ AI trả lời đơn lẻ sang AI hiểu toàn bộ hành trình học tập.

\begin{figure}[H]
\centering
\includegraphics[width=1.0\textwidth]{khao_sat_muc_do_su_dung_ai_trong_hoc_tap.png}
\caption{Khảo sát mức độ sử dụng AI trong học tập}
\label{fig:khao-sat-muc-do-su-dung-ai-trong-hoc-tap}
\end{figure}

\subsubsection{Về kỳ vọng chức năng của AI Tutor}

Khi được hỏi về các chức năng mong muốn, phân tích lỗi code và gợi ý sửa lỗi là hai tính năng được nhắc đến nhiều nhất, tiếp theo là giải thích khái niệm lập trình và gợi ý bài học tiếp theo. Một số phản hồi mở đề cập đến nhu cầu về lộ trình học tập chi tiết, nhắc nhở học tập định kỳ và theo dõi tiến độ theo thời gian, đây đều là những tính năng mà các công cụ AI thông thường không cung cấp. Tập hợp này chính là cơ sở để xác định các chức năng cốt lõi của hệ thống PPAT.

\begin{figure}[H]
\centering
\includegraphics[width=1.0\textwidth]{khao_sat_ky_vong_chuc_nang_cua_ai_tutor.png}
\caption{Khảo sát kỳ vọng chức năng của AI Tutor}
\label{fig:khao-sat-ky-vong-chuc-nang-cua-ai-tutor}
\end{figure}

\subsubsection{Về sự sẵn sàng sử dụng}

Gần 57\% người được hỏi cho biết sẵn sàng sử dụng một hệ thống AI Tutor chuyên biệt, và thêm 36\% trả lời có thể cân nhắc, chỉ có 8\% từ chối. Tỷ lệ muốn tham gia thử nghiệm hệ thống lên đến 93\%. Những con số này cho thấy nhu cầu thực sự tồn tại và người dùng mục tiêu có thái độ cởi mở, tạo điều kiện thuận lợi cho việc triển khai và thu thập phản hồi trong giai đoạn thử nghiệm.

\begin{figure}[H]
\centering
\includegraphics[width=1.0\textwidth]{khao_sat_su_san_sang_su_dung.png}
\caption{Khảo sát sự sẵn sàng sử dụng}
\label{fig:khao-sat-su-san-sang-su-dung}
\end{figure}

\subsection{Sơ đồ BPMN quy trình học tập}

Để làm rõ sự khác biệt giữa cách học lập trình hiện tại và cách học có sự hỗ trợ của hệ thống PPAT, tác giả xây dựng hai sơ đồ quy trình nghiệp vụ theo chuẩn BPMN 2.0, lần lượt mô tả hai luồng hoạt động tương ứng.

\begin{figure}[H]
\centering
\includegraphics[width=1.0\textwidth]{quy_trinh_hoc_lap_trinh_hien_tai.png}
\caption{Quy trình học lập trình hiện tại}
\label{fig:quy-trinh-hoc-lap-trinh-hien-tai}
\end{figure}

Sơ đồ BPMN thứ nhất (Hình 2.11) thể hiện quy trình học lập trình phổ biến mà sinh viên đang thực hiện trước khi có hệ thống hỗ trợ chuyên biệt. Quy trình bắt đầu khi sinh viên tiếp nhận bài học lý thuyết từ giảng viên hoặc tài nguyên tự học như video, tài liệu. Sau đó sinh viên chuyển sang giai đoạn làm bài tập thực hành. Tại đây xuất hiện điểm rẽ nhánh quan trọng: nếu bài tập hoàn thành thành công, sinh viên tiếp tục sang bài học tiếp theo; nếu gặp lỗi hoặc không hiểu, sinh viên phải tự tìm kiếm giải pháp thông qua Google hoặc đặt câu hỏi cho các công cụ AI thông thường như ChatGPT.

Điểm yếu cốt lõi của luồng này nằm ở bước tìm kiếm hỗ trợ. Mỗi phiên tìm kiếm diễn ra hoàn toàn độc lập, không có bất kỳ thông tin nào về lịch sử học tập, trình độ hiện tại hay các lỗi mà sinh viên đó đã từng mắc phải. Sau khi nhận được câu trả lời, sinh viên quay lại làm bài tập nhưng không có cơ chế nào ghi nhận liệu vấn đề đã thực sự được giải quyết triệt để hay chưa. Vòng lặp này có thể lặp đi lặp lại nhiều lần với cùng một dạng lỗi mà không tạo ra sự cải thiện bền vững trong khả năng của người học.

\begin{figure}[H]
\centering
\includegraphics[width=1.0\textwidth]{quy_trinh_hoc_lap_trinh_voi_he_thong_ppat.png}
\caption{Quy trình học lập trình với hệ thống PPAT}
\label{fig:quy-trinh-hoc-lap-trinh-voi-he-thong-ppat}
\end{figure}

Sơ đồ BPMN thứ hai (Hình 2.12) mô tả quy trình học tập khi sinh viên sử dụng hệ thống PPAT. Điểm khởi đầu tương tự: sinh viên tiếp nhận bài học và chuyển sang làm bài tập. Tuy nhiên từ đây, luồng xử lý có sự khác biệt căn bản.

Khi sinh viên gặp khó khăn và đặt câu hỏi cho AI Tutor, hệ thống không xử lý yêu cầu như một câu hỏi đơn lẻ. Trước khi sinh ra phản hồi, hệ thống truy vấn knowledge log của sinh viên đó từ cơ sở dữ liệu, bao gồm trình độ hiện tại, các chủ đề đã nắm vững, các lỗi thường gặp và lịch sử tương tác gần đây. Dựa trên tập thông tin này, AI Tutor xây dựng phản hồi được cá nhân hóa theo đúng ngữ cảnh của người học. Sau mỗi lần tương tác, hệ thống tự động cập nhật lại knowledge log, ghi nhận thêm thông tin mới về hành vi và tiến độ của sinh viên.

Khi sinh viên nộp bài tập, hệ thống chấm qua Code Runner và phân nhánh theo kết quả. Nếu vượt qua, bài học tiếp theo được mở khóa và knowledge log được cập nhật theo hướng tích cực. Nếu thất bại, hệ thống ghi nhận loại lỗi, so sánh với lịch sử lỗi đã có và điều chỉnh gợi ý phù hợp cho lần tiếp theo. Toàn bộ vòng lặp này tạo ra một cơ chế học tập khép kín, trong đó mỗi tương tác đều đóng góp vào việc làm phong phú thêm hồ sơ học tập của sinh viên và nâng cao chất lượng hỗ trợ ở những phiên sau.

So sánh hai sơ đồ cho thấy sự khác biệt không chỉ nằm ở công nghệ sử dụng mà còn ở triết lý thiết kế. Quy trình cũ xử lý từng sự kiện học tập một cách rời rạc, trong khi quy trình mới kết nối tất cả các sự kiện đó thành một hành trình học tập liên tục có trí nhớ và khả năng thích nghi.

\subsection{Phân tích hệ thống hiện có}

Trước khi xây dựng hệ thống PPAT, em tiến hành phân tích các công cụ hỗ trợ học lập trình hiện có mà sinh viên đang sử dụng phổ biến, trong đó ChatGPT là đại diện tiêu biểu nhất. Mục tiêu của phân tích này là xác định rõ những giới hạn của công cụ hiện tại, từ đó làm cơ sở lý giải sự cần thiết của một hệ thống chuyên biệt hơn.

ChatGPT và các công cụ AI tổng quát tương tự đang được phần lớn sinh viên sử dụng như một trợ lý học tập không chính thức. Kết quả khảo sát ở mục trước cho thấy hơn 53\% sinh viên chọn AI làm điểm dừng đầu tiên khi gặp lỗi code. Điều này phản ánh một thực tế rõ ràng: các công cụ này có khả năng trả lời nhanh, phủ rộng nhiều chủ đề và dễ tiếp cận. Tuy nhiên, chính vì được thiết kế cho mục đích tổng quát nên chúng bộc lộ nhiều hạn chế khi được dùng trong ngữ cảnh học tập có hệ thống.

Hạn chế lớn nhất là không có khả năng ghi nhớ lịch sử học tập xuyên suốt nhiều phiên làm việc. Mỗi cuộc trò chuyện với ChatGPT bắt đầu lại từ đầu, không có bất kỳ thông tin nào về người dùng là ai, đang học đến đâu, hay đã gặp khó khăn gì trước đó. Điều này khiến mỗi phản hồi đều mang tính chung chung, không được điều chỉnh theo trình độ thực tế của người học. Một sinh viên mới bắt đầu học lập trình và một sinh viên năm ba có thể nhận được câu trả lời gần như giống nhau cho cùng một câu hỏi, dù nhu cầu và nền tảng kiến thức của hai người hoàn toàn khác biệt.

Hạn chế thứ hai liên quan đến chiến lược dạy học. ChatGPT có xu hướng cung cấp đáp án trực tiếp thay vì dẫn dắt người học tự tìm ra giải pháp. Cách tiếp cận này tuy nhanh nhưng không giúp sinh viên hình thành tư duy giải quyết vấn đề, vốn là kỹ năng cốt lõi mà bất kỳ lập trình viên nào cũng cần có. Sinh viên nhận được code hoạt động nhưng không hiểu tại sao code cũ của mình sai và lần sau gặp vấn đề tương tự vẫn phải hỏi lại từ đầu.

Hạn chế thứ ba là không có cơ chế theo dõi tiến độ và phát hiện điểm yếu. ChatGPT không biết sinh viên đó đã hoàn thành bao nhiêu bài tập, hay mắc lỗi ở dạng bài nào nhiều nhất. Do đó không thể chủ động đề xuất nội dung ôn tập, điều chỉnh độ khó hay cảnh báo khi phát hiện một lỗ hổng kiến thức đang lặp lại. Người học hoàn toàn phải tự điều hướng quá trình học của mình mà không có bất kỳ phản hồi có cấu trúc nào từ hệ thống.

Ngoài ra, ChatGPT không tích hợp môi trường thực thi code. Sinh viên phải chạy code trên môi trường riêng rồi sao chép kết quả hoặc thông báo lỗi vào chat, tạo ra một bước thủ công làm gián đoạn luồng học tập. Trong khi đó, việc thực thi và kiểm tra code ngay trong cùng một giao diện là yếu tố quan trọng để tạo ra vòng phản hồi nhanh, giúp sinh viên học qua thực hành hiệu quả hơn.

\section{Xác định Actor}

Hệ thống PPAT được thiết kế để phục vụ hai nhóm người dùng với vai trò và nhu cầu khác nhau. Việc xác định rõ các actor tham gia vào hệ thống là bước quan trọng để đảm bảo các chức năng được xây dựng đúng đối tượng và đáp ứng đúng nhu cầu thực tế.

Actor thứ nhất là Sinh viên, người dùng trung tâm và quan trọng nhất của hệ thống. Sinh viên tương tác trực tiếp với toàn bộ các chức năng cốt lõi bao gồm xem bài học, làm bài tập, chat với AI Tutor và theo dõi tiến độ học tập của bản thân. Để truy cập bất kỳ chức năng nào, sinh viên bắt buộc phải đăng ký tài khoản và đăng nhập vào hệ thống. Mọi cơ chế cá nhân hóa trong hệ thống đều hướng đến việc phục vụ tốt nhất cho nhóm actor này.

Actor thứ hai là Admin đóng vai trò quản trị toàn bộ hệ thống. Admin có quyền quản lý tài khoản người dùng, tạo mới, chỉnh sửa và xóa bài học, bài tập cùng các test case tương ứng. Khác với sinh viên chỉ có quyền đọc và tương tác với nội dung, Admin có toàn quyền kiểm soát kho nội dung học tập. Admin cũng cần đăng nhập vào hệ thống trước khi thực hiện bất kỳ thao tác quản trị nào.

\section{Yêu cầu chức năng}

Dựa trên kết quả khảo sát thực tế và phân tích các actor tham gia hệ thống, em xác định các nhóm yêu cầu chức năng chính:

\subsubsection{Quản lý người dùng}

\begin{itemize}
    \item Đăng ký tài khoản bằng email và mật khẩu.
    \item Đăng nhập, xác thực bằng JWT với thời hạn 24 giờ và refresh token 7 ngày.
    \item Cập nhật thông tin cá nhân, mật khẩu, trình độ ban đầu.
\end{itemize}

\subsubsection{Quản lý bài học}

\begin{itemize}
    \item Xem danh sách bài học, lọc theo ngôn ngữ lập trình, chủ đề, mức độ khó.
    \item Xem chi tiết bài học với nội dung lý thuyết và ví dụ code có syntax highlighting.
    \item Theo dõi trạng thái hoàn thành của từng bài học.
\end{itemize}

\subsubsection{Quản lý bài tập}

\begin{itemize}
    \item Xem đề bài kèm starter code tương ứng.
    \item Viết và chỉnh sửa code trực tiếp trong Monaco Editor tích hợp.
    \item Nộp bài và kiểm tra kết quả tự động qua test cases.
    \item Nhận gợi ý sửa lỗi theo mức độ tương ứng với số lần thất bại liên tiếp.
\end{itemize}

\subsubsection{Tương tác chat với AI Tutor}

\begin{itemize}
    \item Gửi câu hỏi và nhận phản hồi cá nhân hóa từ GPT-4o theo trình độ và lịch sử học tập.
    \item Nhận phản hồi theo cơ chế streaming token-by-token, không chờ toàn bộ nội dung.
    \item Ghi nhớ Context, lưu lịch sử hội thoại vĩnh viễn, có thể xem lại bất kỳ lúc nào.
    \item Tạo nhiều phiên hội thoại độc lập cho các chủ đề khác nhau.
\end{itemize}

\subsubsection{Theo dõi tiến độ học tập}

\begin{itemize}
    \item Xem tổng quan tiến độ: số bài đã hoàn thành, chuỗi ngày học liên tiếp.
    \item Xem biểu đồ radar thể hiện điểm mạnh và điểm yếu theo từng chủ đề.
    \item Xem danh sách các lỗi thường gặp được phân tích từ lịch sử nộp bài.
    \item Xem đề xuất bài học tiếp theo dựa trên knowledge log hiện tại.
\end{itemize}

\subsubsection{Quản lý nội dung (Admin)}

\begin{itemize}
    \item Tạo mới, chỉnh sửa và xóa bài học, bài tập và test cases.
    \item Phân loại nội dung theo ngôn ngữ, chủ đề và mức độ khó.
\end{itemize}

\subsubsection{Cập nhật Knowledge Log (Hệ thống tự động)}

\begin{itemize}
    \item Tự động cập nhật hồ sơ học tập sau mỗi lần tương tác chat hoặc nộp bài.
    \item Ghi nhận loại lỗi, tần suất mắc lỗi và tốc độ tiến bộ của từng sinh viên.
    \item Điều chỉnh lộ trình học tập gợi ý dựa trên dữ liệu tích lũy theo thời gian.
\end{itemize}


{
\renewcommand{\arraystretch}{1.5}
\begin{tabularx}{\textwidth}{|p{2.0cm}|p{4cm}|X|}
\caption{Danh sách yêu cầu chức năng}
\label{tab:danh-sach-yeu-cau-chuc-nang}
\\
\hline
\textbf{Mã} & \textbf{Tên yêu cầu} & \textbf{Mô tả ngắn} \\
\hline
FR01 & Quản lý người dùng & Đăng ký, đăng nhập JWT, cập nhật profile \\
\hline
FR02 & Quản lý bài học & Danh sách bài học, lọc theo ngôn ngữ/chủ đề/độ khó \\
\hline
FR03 & Quản lý bài tập & Xem đề, code editor, submit, chạy test case, xem gợi ý \\
\hline
FR04 & Chat với AI Tutor & Chat streaming SSE, phản hồi cá nhân hóa, lưu lịch sử \\
\hline
FR05 & Theo dõi tiến độ & Dashboard: streak, bài hoàn thành, điểm mạnh/yếu \\
\hline
FR06 & Quản lý nội dung & CRUD bài học, bài tập, test cases, import/export \\
\hline
FR07 & Cập nhật Knowledge Log & Tự động sau mỗi interaction, trigger điều chỉnh lộ trình \\
\hline
\end{tabularx}
}

\section{Yêu cầu phi chức năng}

Bên cạnh các yêu cầu chức năng, hệ thống PPAT còn phải đáp ứng một tập hợp các yêu cầu phi chức năng liên quan đến hiệu năng, bảo mật, khả dụng và trải nghiệm người dùng.

\subsubsection{Hiệu năng}

\begin{itemize}
    \item Thời gian phản hồi của các API không gọi LLM không vượt quá 200 mili giây trong điều kiện tải bình thường.
    \item Hệ thống xử lý đồng thời tối thiểu 50 người dùng trong môi trường triển khai phát triển.
    \item Chức năng chat AI sử dụng cơ chế streaming, người dùng nhận được token đầu tiên ngay lập tức thay vì chờ toàn bộ phản hồi.
\end{itemize}

\subsubsection{Bảo mật}

\begin{itemize}
    \item Xác thực người dùng bằng JWT với thời hạn 24 giờ, kết hợp refresh token có hiệu lực 7 ngày.
    \item Mật khẩu được mã hóa bằng bcrypt trước khi lưu vào cơ sở dữ liệu.
    \item Code của sinh viên được thực thi trong Docker container cô lập, giới hạn thời gian chạy 10 giây, không có quyền truy cập mạng và không thể ghi ra filesystem của máy chủ.
    \item Các endpoint nhạy cảm được bảo vệ bằng rate limiting để ngăn chặn tấn công brute force.
\end{itemize}

\subsubsection{Tính ổn định và tin cậy}

\begin{itemize}
    \item Hệ thống duy trì uptime 99\% trong khung giờ học tập từ 7 giờ sáng đến 10 giờ tối.
    \item Dữ liệu PostgreSQL được sao lưu tự động hàng ngày.
    \item Redis được cấu hình chế độ AOF để đảm bảo không mất dữ liệu session khi khởi động lại.
\end{itemize}

\subsubsection{Khả năng mở rộng}

\begin{itemize}
    \item Kiến trúc SOA cho phép thêm service mới như analytics hoặc notification độc lập mà không ảnh hưởng đến các thành phần hiện tại.
    \item Code Runner có thể được tách thành microservice riêng biệt hoặc mở rộng hỗ trợ thêm ngôn ngữ lập trình khi cần.
\end{itemize}

\subsubsection{Khả năng sử dụng}

\begin{itemize}
    \item Giao diện trực quan, thao tác nhất quán giữa các màn hình.
    \item Monaco Editor tích hợp sẵn syntax highlighting và hỗ trợ nhiều ngôn ngữ lập trình.
    \item Phản hồi streaming tạo trải nghiệm tự nhiên, giảm cảm giác chờ đợi khi tương tác với AI.
\end{itemize}

\subsubsection{Khả năng tương thích}

\begin{itemize}
    \item Hoạt động tốt trên các trình duyệt phổ biến bao gồm Chrome, Edge và Firefox.
    \item Giao diện responsive, hỗ trợ sử dụng trên cả máy tính và thiết bị di động.
\end{itemize}

\section{Mô hình nghiệp vụ}

\subsection{Biểu đồ Use Case}

\begin{itemize}
    \item Biểu đồ UC tổng quát
\end{itemize}

\begin{figure}[H]
\centering
\includegraphics[width=1.0\textwidth]{bieu_do_uc_tong_quat.png}
\caption{Biểu đồ UC tổng quát}
\label{fig:bieu-do-uc-tong-quat}
\end{figure}

\begin{itemize}
    \item Biểu đồ UC phân rã
    \item Biểu đồ phân rã cho tác nhân Admin
\end{itemize}

\begin{figure}[H]
\centering
\includegraphics[width=1.0\textwidth]{bieu_do_phan_ra_cho_tac_nhan_admin.png}
\caption{Biểu đồ phân rã cho tác nhân Admin}
\label{fig:bieu-do-phan-ra-cho-tac-nhan-admin}
\end{figure}

\begin{itemize}
    \item 1. Biểu đồ phân rã cho tác nhân Sinh viên
\end{itemize}

\begin{figure}[H]
\centering
\includegraphics[width=1.0\textwidth]{bieu_do_phan_ra_cho_tac_nhan_sinh_vien.png}
\caption{Biểu đồ phân rã cho tác nhân Sinh viên}
\label{fig:bieu-do-phan-ra-cho-tac-nhan-sinh-vien}
\end{figure}

\subsection{Đặc tả Use Case}

\subsubsection{Đặc tả chức năng Đăng ký (Sinh viên)}


{
\renewcommand{\arraystretch}{1.5}
\begin{tabularx}{\textwidth}{|p{3.5cm}|X|}
\caption{Bảng đặc tả chức năng Đăng ký (Sinh viên)}
\label{tab:bang-dac-ta-chuc-nang-dang-ky-sinh-vien}
\\
\hline
\textbf{Mã Use Case} & UC01 \\
\hline
\textbf{Tên Use Case} & Đăng ký tài khoản \\
\hline
\textbf{Tác nhân chính} & Sinh viên \\
\hline
\textbf{Mô tả} & Người dùng (Sinh viên) thực hiện tạo tài khoản mới trên hệ thống để bắt đầu quá trình học tập. \\
\hline
\textbf{Tiền điều kiện} & - Người dùng chưa có tài khoản trong hệ thống. \newline - Hệ thống kết nối ổn định với CSDL. \\
\hline
\textbf{Hậu điều kiện} & - Một bản ghi người dùng mới được tạo trong CSDL với vai trò là student. \newline - Hệ thống cấp mã thông báo (Access Token) và tự động đăng nhập cho người dùng. \\
\hline
\textbf{Luồng chính} & 1. Người dùng chọn chức năng "Đăng ký" từ giao diện. \newline 2. Hệ thống hiển thị form đăng ký bao gồm các trường: Email, Mật khẩu, Tên hiển thị, Trình độ (Beginner/Intermediate/Advanced). \newline 3. Người dùng nhập đầy đủ và chính xác các thông tin yêu cầu. \newline 4. Người dùng nhấn nút "Đăng ký". \newline 5. Hệ thống thực hiện kiểm tra dữ liệu: \newline 5.1. Kiểm tra các trường bắt buộc không được để trống. \newline 5.2. Kiểm tra định dạng Email hợp lệ. \newline 5.3. Kiểm tra tính duy nhất của Email trong CSDL. \newline 6. Hệ thống thực hiện mã hóa mật khẩu (bằng bcrypt). \newline 7. Hệ thống lưu thông tin người dùng mới vào CSDL. \newline 8. Hệ thống tạo Access Token và Refresh Token. \newline 9. Hệ thống hiển thị thông báo "Đăng ký thành công" và chuyển hướng người dùng vào trang chủ. \\
\hline
\textbf{Luồng phụ} & - 5.1.a. Nếu thông tin bị thiếu: Hệ thống hiển thị thông báo "Vui lòng điền đầy đủ các trường bắt buộc". \newline - 5.2.a. Nếu Email sai định dạng: Hệ thống hiển thị cảnh báo "Email không hợp lệ". \newline - 5.3.a. Nếu Email đã tồn tại: Hệ thống hiển thị lỗi "Email đã được sử dụng". \\
\hline
\textbf{Ngoại lệ} & - 7.a. Nếu quá trình lưu vào CSDL thất bại (lỗi server): Hệ thống thông báo "Đăng ký thất bại, vui lòng thử lại sau". \newline - 7.b. Mất kết nối CSDL: Hệ thống hiển thị lỗi kết nối máy chủ. \\
\hline
\end{tabularx}
}

\subsubsection{Đặc tả chức năng Đăng nhập}


{
\renewcommand{\arraystretch}{1.5}
\begin{tabularx}{\textwidth}{|p{3.5cm}|X|}
\caption{Bảng đặc tả chức năng Đăng nhập (Sinh viên, Admin)}
\label{tab:bang-dac-ta-chuc-nang-dang-nhap-sinh-vien-admin}
\\
\hline
\textbf{Mã Use Case} & UC01 \\
\hline
\textbf{Tên Use Case} & Đăng nhập hệ thống \\
\hline
\textbf{Tác nhân chính} & Sinh viên, Admin \\
\hline
\textbf{Mô tả} & Người dùng thực hiện đăng nhập vào hệ thống để truy cập vào các tính năng theo phân quyền của mình. \\
\hline
\textbf{Tiền điều kiện} & - Người dùng đã có tài khoản hợp lệ trong hệ thống. \newline - Hệ thống đang kết nối ổn định với CSDL. \\
\hline
\textbf{Hậu điều kiện} & - Người dùng truy cập thành công vào hệ thống. \newline - Hệ thống trả về thông tin hồ sơ và mã xác thực (Access Token, Refresh Token). \\
\hline
\textbf{Luồng chính} & 1. Người dùng truy cập vào trang Đăng nhập. \newline 2. Hệ thống hiển thị form đăng nhập bao gồm: Email và Mật khẩu. \newline 3. Người dùng nhập Email và Mật khẩu sau đó nhấn nút "Đăng nhập". \newline 4. Hệ thống kiểm tra thông tin đăng nhập: \newline 4.1. Truy vấn thông tin người dùng theo Email. \newline 4.2. So sánh mật khẩu đã nhập với bản mã hóa trong CSDL (sử dụng bcrypt.compare). \newline 5. Nếu thông tin khớp: \newline 5.1. Hệ thống tạo phiên làm việc mới (Access Token 24h, Refresh Token 7d). \newline 5.2. Hệ thống xác định vai trò (Role: admin/student) và các quyền hạn tương ứng. \newline 5.3. Hệ thống trả về thông báo "Đăng nhập thành công". \newline 6. Hệ thống chuyển hướng người dùng về trang chủ hoặc trang Dashboard theo vai trò. \\
\hline
\textbf{Luồng phụ} & - 5.a. Nếu Email hoặc Mật khẩu không chính xác: Hệ thống hiển thị thông báo "Sai email hoặc mật khẩu" và yêu cầu nhập lại. \newline - 5.b. Nếu phát hiện Refresh Token cũ còn hiệu lực: Hệ thống thực hiện thu hồi và cấp mã mới để đảm bảo bảo mật. \\
\hline
\textbf{Ngoại lệ} & - 4.1.a. Nếu CSDL bị lỗi: Hệ thống hiển thị thông báo lỗi hệ thống chung. \newline - 4.1.b. Mất kết nối mạng: Hiển thị lỗi kết nối máy chủ. \\
\hline
\end{tabularx}
}

\subsubsection{Đặc tả chức năng Quản lý thông tin cá nhân}


{
\renewcommand{\arraystretch}{1.5}
\begin{tabularx}{\textwidth}{|p{3.5cm}|X|}
\caption{Bảng đặc tả chức năng Quản lý thông tin cá nhân}
\label{tab:bang-dac-ta-chuc-nang-quan-ly-thong-tin-ca-nhan}
\\
\hline
\textbf{Mã Use Case} & UC01 \\
\hline
\textbf{Tên Use Case} & Quản lý thông tin cá nhân \\
\hline
\textbf{Tác nhân chính} & Sinh viên, Admin \\
\hline
\textbf{Mô tả} & Cho phép người dùng tự cập nhật thông tin cá nhân của mình bao gồm tên hiển thị, mật khẩu và trình độ học tập hiện tại. \\
\hline
\textbf{Tiền điều kiện} & Người dùng đã đăng nhập thành công vào hệ thống. \\
\hline
\textbf{Hậu điều kiện} & Thông tin cá nhân của sinh viên, Admin được cập nhật thành công trong CSDL. \\
\hline
\textbf{Luồng chính} & 1. Người dùng chọn chức năng "Thông tin tài khoản". \newline 2. Hệ thống hiển thị thông tin hiện tại: Tên hiển thị, Trình độ hiện tại. \newline 3. Người dùng thực hiện các thay đổi mong muốn: \newline 3.1. Thay đổi tên hiển thị. \newline 3.2. Thay đổi trình độ học tập (Beginner/Intermediate/Advanced). \newline 3.3. Thay đổi mật khẩu (yêu cầu nhập mật khẩu cũ và mật khẩu mới). \newline 4. Người dùng nhấn "Lưu thay đổi". \newline 5. Hệ thống kiểm tra tính hợp lệ của dữ liệu: \newline 5.1. Kiểm tra mật khẩu cũ có khớp với dữ liệu trong CSDL hay không. \newline 5.2. Kiểm tra độ mạnh của mật khẩu mới (nếu có thay đổi). \newline 6. Hệ thống cập nhật thông tin mới vào CSDL. \newline 7. Hệ thống hiển thị thông báo "Cập nhật thông tin thành công". \\
\hline
\textbf{Luồng phụ} & - 5.1.a. Nếu mật khẩu cũ không chính xác: Hệ thống báo lỗi "Mật khẩu cũ không chính xác" và yêu cầu nhập lại. \newline - 5.2.a. Nếu thông tin nhập vào không hợp lệ (ví dụ tên quá ngắn): Hệ thống hiển thị cảnh báo tương ứng. \\
\hline
\textbf{Ngoại lệ} & - 6.a. Nếu quá trình cập nhật gặp lỗi CSDL: Hệ thống thông báo lỗi hệ thống và yêu cầu thử lại sau. \newline - 6.b. Mất kết nối mạng: Hiển thị lỗi kết nối máy chủ. \\
\hline
\end{tabularx}
}

\subsubsection{Đặc tả chức năng Quản lý người dùng}


{
\renewcommand{\arraystretch}{1.5}
\begin{tabularx}{\textwidth}{|p{3.5cm}|X|}
\caption{Bảng đặc tả chức năng Quản lý người dùng}
\label{tab:bang-dac-ta-chuc-nang-quan-ly-nguoi-dung}
\\
\hline
\textbf{Mã Use Case} & UC06 \\
\hline
\textbf{Tên Use Case} & Quản lý người dùng \\
\hline
\textbf{Tác nhân chính} & Admin \\
\hline
\textbf{Mô tả} & Admin thực hiện quản lý toàn diện (thêm, sửa, xóa) các tài khoản người dùng khác trong hệ thống. \\
\hline
\textbf{Tiền điều kiện} & Admin đã đăng nhập thành công với quyền quản trị viên. \\
\hline
\textbf{Hậu điều kiện} & Các thay đổi về thông tin người dùng khác được lưu thành công vào CSDL. \\
\hline
\textbf{Luồng chính} & 1. Thêm người dùng mới: \newline 1.1. Admin chọn "Thêm người dùng". \newline 1.2. Nhập Email, Tên, Vai trò (Admin/Student), Mật khẩu và Trình độ. \newline 1.3. Hệ thống kiểm tra tính duy nhất của Email và lưu vào CSDL. \newline 2. Sửa thông tin người dùng khác: \newline 2.1. Admin chọn người dùng cần sửa từ danh sách quản lý. \newline 2.2. Hệ thống hiển thị form chỉnh sửa. \newline 2.3. Admin thay đổi thông tin (trừ Email) và nhấn "Cập nhật". \newline 2.4. Hệ thống ghi đè thông tin mới vào CSDL. \newline 3. Xóa tài khoản người dùng: \newline 3.1. Admin chọn tài khoản cần xóa và bấm "Xóa". \newline 3.2. Hệ thống hiển thị hộp thoại xác nhận xóa. \newline 3.3. Admin xác nhận "Đồng ý". \newline 3.4. Hệ thống xóa vĩnh viễn dữ liệu người dùng và các thông tin liên quan, báo "Xóa thành công". \\
\hline
\textbf{Luồng phụ} & - 1.2.1. Email trùng: Báo lỗi "Email đã được đăng ký cho tài khoản khác". \newline - 3.3.1. Admin chọn "Hủy": Đóng hộp thoại, không thay đổi dữ liệu. \\
\hline
\textbf{Ngoại lệ} & - 2.4.a. Lỗi hệ thống khi cập nhật: Thông báo "Cập nhật thất bại, vui lòng thử lại". \newline - 3.4.a. Mất kết nối CSDL: Hiển thị lỗi kết nối máy chủ. \\
\hline
\end{tabularx}
}

\subsubsection{Đặc tả chức năng Quản lý bài học}


{
\renewcommand{\arraystretch}{1.5}
\begin{tabularx}{\textwidth}{|p{3.5cm}|X|}
\caption{Bảng đặc tả chức năng Quản lý bài học}
\label{tab:bang-dac-ta-chuc-nang-quan-ly-bai-hoc}
\\
\hline
\textbf{Mã Use Case} & UC06 \\
\hline
\textbf{Tên Use Case} & Quản lý bài học \\
\hline
\textbf{Tác nhân chính} & Admin \\
\hline
\textbf{Mô tả} & Usecase này cho phép Admin thực hiện các thao tác thêm, sửa, xóa các bài học trong hệ thống để cung cấp nội dung học tập cho người dùng. \\
\hline
\textbf{Tiền điều kiện} & Admin đã đăng nhập thành công vào hệ thống với quyền quản trị viên. \\
\hline
\textbf{Hậu điều kiện} & Thông tin bài học được cập nhật thành công trong CSDL. \\
\hline
\textbf{Luồng chính} & 1. Thêm Bài học: \newline 1.1. Admin chọn chức năng "Thêm bài học". \newline 1.2. Hệ thống hiển thị biểu mẫu bao gồm: Tiêu đề, Ngôn ngữ, Chủ đề, Độ khó (Dễ/Trung bình/Khó) và Nội dung (Markdown). \newline 1.3. Admin nhập đầy đủ thông tin và chọn "Lưu". \newline 1.4. Hệ thống kiểm tra dữ liệu: \newline 1.4.1. Nếu hợp lệ: Lưu thông tin bài học mới vào CSDL và báo "Thêm bài học thành công". \newline 1.4.2. Nếu không hợp lệ: Hiển thị thông báo lỗi. \newline 2. Sửa Bài học: \newline 2.1. Admin chọn bài học cần chỉnh sửa từ danh sách. \newline 2.2. Hệ thống hiển thị thông tin bài học hiện tại. \newline 2.3. Admin thực hiện thay đổi và chọn "Cập nhật". \newline 2.4. Hệ thống kiểm tra và ghi đè dữ liệu mới vào CSDL, báo "Cập nhật thành công". \newline 3. Xóa Bài học: \newline 3.1. Admin chọn bài học và yêu cầu xóa. \newline 3.2. Hệ thống hiển thị hộp thoại xác nhận. \newline 3.3. Admin chọn "Đồng ý": \newline 3.3.1. Hệ thống xóa bài học và các dữ liệu liên quan (tập bài tập, tiến độ) khỏi hệ thống. \newline 3.3.2. Hiển thị thông báo "Xóa bài học thành công". \\
\hline
\textbf{Luồng phụ} & - 1.3.1. Thiếu thông tin bắt buộc: Hệ thống yêu cầu người dùng điền đầy đủ. \newline - 3.3.2. Admin chọn "Hủy": Đóng hộp thoại, không thay đổi dữ liệu. \\
\hline
\textbf{Ngoại lệ} & - 1.4.a. Lỗi kết nối CSDL: Hệ thống hiển thị thông báo lỗi kết nối. \newline - 3.3.1.a. Xóa thất bại do lỗi ràng buộc: Thông báo lỗi hệ thống. \\
\hline
\end{tabularx}
}

\subsubsection{Đặc tả chức năng Quản lý bài tập}


{
\renewcommand{\arraystretch}{1.5}
\begin{tabularx}{\textwidth}{|p{3.5cm}|X|}
\caption{Bảng đặc tả chức năng Quản lý bài tập}
\label{tab:bang-dac-ta-chuc-nang-quan-ly-bai-tap}
\\
\hline
\textbf{Mã Use Case} & UC06 \\
\hline
\textbf{Tên Use Case} & Quản lý bài tập \\
\hline
\textbf{Tác nhân chính} & Admin \\
\hline
\textbf{Mô tả} & Usecase này cho phép Admin thiết kế và quản lý các bài tập đi kèm với bài học, bao gồm các thông tin về mã nguồn khởi tạo, gợi ý và kết quả mong muốn. \\
\hline
\textbf{Tiền điều kiện} & Admin đã đăng nhập thành công vào hệ thống. \\
\hline
\textbf{Hậu điều kiện} & Thông tin tập bài tập được lưu trữ hoặc cập nhật thành công trong CSDL. \\
\hline
\textbf{Luồng chính} & 1. Thêm Bài Tập: \newline 1.1. Admin chọn bài học muốn thêm bài tập và chọn "Thêm bài tập". \newline 1.2. Hệ thống hiển thị biểu mẫu: Tiêu đề, Mô tả, Mã nguồn khởi tạo, Kết quả mong muốn, Gợi ý, Bài học tiên quyết. \newline 1.3. Admin nhập các thông tin chi tiết và chọn "Lưu". \newline 1.4. Hệ thống kiểm tra và lưu vào CSDL, báo "Thêm bài tập thành công". \newline 2. Sửa Bài Tập: \newline 2.1. Admin chọn bài tập cần chỉnh sửa. \newline 2.2. Hệ thống tải dữ liệu hiện tại lên biểu mẫu. \newline 2.3. Admin thay đổi nội dung (ví dụ gợi ý hoặc mã khởi tạo) và nhấn "Cập nhật". \newline 2.4. Hệ thống lưu các thay đổi và thông báo thành công. \newline 3. Xóa Bài Tập: \newline 3.1. Admin chọn bài tập cần xóa. \newline 3.2. Hệ thống yêu cầu xác nhận. \newline 3.3. Admin xác nhận "Xóa vĩnh viễn": \newline 3.3.1. Hệ thống xóa bài tập và các bản nộp bài (submission) liên quan của người dùng. \newline 3.3.2. Thông báo xóa thành công. \\
\hline
\textbf{Luồng phụ} & - 1.2.1. Không chọn bài học liên kết: Hệ thống yêu cầu chọn bài học để gán bài tập. \newline - 2.2.1. Thông tin nhập sai định dạng: Hệ thống hiển thị cảnh báo. \\
\hline
\textbf{Ngoại lệ} & - 1.4.a. Lỗi hệ thống khi lưu: Hiển thị thông báo thử lại sau. \newline - 3.3.1.a. Mất kết nối CSDL: Hiển thị lỗi kết nối. \\
\hline
\end{tabularx}
}

\subsubsection{Đặc tả chức năng Quản lý test cases}


{
\renewcommand{\arraystretch}{1.5}
\begin{tabularx}{\textwidth}{|p{3.5cm}|X|}
\caption{Bảng đặc tả chức năng Quản lý test cases}
\label{tab:bang-dac-ta-chuc-nang-quan-ly-test-cases}
\\
\hline
\textbf{Mã Use Case} & UC06 \\
\hline
\textbf{Tên Use Case} & Quản lý test cases \\
\hline
\textbf{Tác nhân chính} & Admin \\
\hline
\textbf{Mô tả} & Cho phép Admin thiết lập các bộ dữ liệu kiểm thử (đầu vào và đầu ra mong muốn) cho từng bài tập để hệ thống tự động chấm điểm bài nộp của người dùng. \\
\hline
\textbf{Tiền điều kiện} & Admin đã đăng nhập và bài tập tương ứng đã tồn tại trong hệ thống. \\
\hline
\textbf{Hậu điều kiện} & Danh sách test cases được cập nhật và áp dụng cho quá trình kiểm tra mã nguồn của người dùng. \\
\hline
\textbf{Luồng chính} & 1. Thêm bộ kiểm thử: \newline 1.1. Admin chọn một bài tập cụ thể và mở mục "Cấu hình Test Cases". \newline 1.2. Admin nhấn "Thêm bộ kiểm thử mới". \newline 1.3. Nhập "Đầu vào (Input)" và "Đầu ra mong muốn (Expected Output)". \newline 1.4. Admin chọn "Lưu". \newline 1.5. Hệ thống thêm cặp dữ liệu vào danh sách kiểm thử của bài tập đó trong CSDL. \newline 2. Sửa bộ kiểm thử: \newline 2.1. Admin chọn test case trong danh sách của bài tập. \newline 2.2. Admin cập nhật giá trị đầu vào hoặc đầu ra. \newline 2.3. Hệ thống ghi lại thay đổi. \newline 3. Xóa bộ kiểm thử: \newline 3.1. Admin nhấn biểu tượng xóa tại một dòng test case. \newline 3.2. Hệ thống thực hiện xóa ngay lập tức khỏi bộ nhớ bản ghi bài tập. \newline 3.3. Hiển thị thông báo "Đã loại bỏ test case". \\
\hline
\textbf{Luồng phụ} & - 1.3.1. Định dạng dữ liệu không khớp: Hệ thống báo lỗi nếu kiểu dữ liệu input/output không hợp lệ cho ngôn ngữ lập trình của bài học. \newline - 1.5.1. Dữ liệu rỗng: Hệ thống yêu cầu nhập đầy đủ trước khi lưu. \\
\hline
\textbf{Ngoại lệ} & - 1.5.a. Lỗi cập nhật CSDL: Thông báo "Lưu cấu hình thất bại". \newline - 3.2.a. Xóa test case cuối cùng: Hệ thống cảnh báo bài tập cần ít nhất một bộ kiểm thử để hoạt động. \\
\hline
\end{tabularx}
}

\subsubsection{Đặc tả chức năng Xem \& Lọc bài học}


{
\renewcommand{\arraystretch}{1.5}
\begin{tabularx}{\textwidth}{|p{3.5cm}|X|}
\caption{Bảng đặc tả chức năng Xem \& Lọc bài học}
\label{tab:bang-dac-ta-chuc-nang-xem-loc-bai-hoc}
\\
\hline
\textbf{Mã Use Case} & UC02 \\
\hline
\textbf{Tên Use Case} & Xem \& Lọc bài học \\
\hline
\textbf{Tác nhân chính} & Sinh viên \\
\hline
\textbf{Mô tả} & Sinh viên truy cập vào thư viện bài học để xem danh sách các bài học, lọc theo tiêu chí ngôn ngữ lập trình và xem chi tiết nội dung lý thuyết của từng bài. \\
\hline
\textbf{Tiền điều kiện} & Sinh viên đã đăng nhập thành công. \\
\hline
\textbf{Hậu điều kiện} & Sinh viên nắm bắt được lộ trình học tập và nội dung lý thuyết của bài học đã chọn. \\
\hline
\textbf{Luồng chính} & 1. Sinh viên chọn mục "Danh sách bài học" trên thanh bên. \newline 2. Hệ thống hiển thị: \newline - Bảng thống kê cá nhân (Số bài đã hoàn thành, số bài tập đã làm, trình độ hiện tại, tỉ lệ hoàn thành). \newline - Danh sách các thẻ bài học. \newline 3. Sinh viên sử dụng bộ lọc ngôn ngữ (Python, JavaScript, C++, Java) hoặc chủ đề/độ khó. \newline 4. Hệ thống cập nhật danh sách bài học khớp với tiêu chí lọc. \newline 5. Sinh viên nhấn vào một thẻ bài học cụ thể. \newline 6. Hệ thống hiển thị trang chi tiết bài học: \newline - Nội dung lý thuyết (định dạng Markdown). \newline - Danh sách các bài tập thực hành liên quan (kèm trạng thái đã hoàn thành hay chưa). \\
\hline
\textbf{Luồng phụ} & - 4.a. Không có bài học nào khớp: Hệ thống hiển thị thông báo "Chưa có bài học nào trong danh mục này". \newline - 6.a. Bài học chưa có bài tập: Hiển thị thông báo "Đang cập nhật bài tập". \\
\hline
\textbf{Ngoại lệ} & - 2.a. Lỗi tải dữ liệu: Hệ thống hiển thị các khung trống (skeleton) và thông báo lỗi kết nối. \\
\hline
\end{tabularx}
}

\subsubsection{Đặc tả chức năng Làm bài tập}


{
\renewcommand{\arraystretch}{1.5}
\begin{tabularx}{\textwidth}{|p{3.5cm}|X|}
\caption{Bảng đặc tả chức năng Làm bài tập}
\label{tab:bang-dac-ta-chuc-nang-lam-bai-tap}
\\
\hline
\textbf{Mã Use Case} & UC03 \\
\hline
\textbf{Tên Use Case} & Làm bài tập \\
\hline
\textbf{Tác nhân chính} & Sinh viên \\
\hline
\textbf{Mô tả} & Sinh viên thực hiện giải quyết các thử thách lập trình thông qua trình soạn thảo code tích hợp, chạy kiểm thử và nhận phản hồi tức thì từ hệ thống. \\
\hline
\textbf{Tiền điều kiện} & Sinh viên đã đăng nhập và chọn một bài tập cụ thể. \\
\hline
\textbf{Hậu điều kiện} & Kết quả làm bài và tiến trình được lưu lại vào hệ thống. \\
\hline
\textbf{Luồng chính} & 1. Sinh viên truy cập vào trang chi tiết một bài tập. \newline 2. Hệ thống hiển thị giao diện gồm: \newline - Cột trái: Tiêu đề, mô tả yêu cầu bài tập, kết quả mong đợi và danh sách bộ kiểm thử (test cases). \newline - Cột phải: Trình soạn thảo mã nguồn (đã có sẵn mã khởi tạo) và nút "Nộp bài", bên dưới là Terminal hiển thị kết quả. \newline 3. Sinh viên viết mã nguồn giải quyết vấn đề vào trình soạn thảo. \newline 4. Sinh viên nhấn nút "Nộp bài". \newline 5. Hệ thống thực hiện: \newline 5.1. Gửi mã nguồn về máy chủ để thực thi. \newline 5.2. Chạy mã nguồn đối chiếu với các bộ kiểm thử. \newline 6. Hệ thống hiển thị kết quả dưới terminal: \newline - Kết quả thực thi code. \newline - Thời gian thực thi. \newline - Chỉ số bộ kiểm thử bị sai (nếu có). \newline - Thông báo lỗi chi tiết từ trình biên dịch/thông dịch (nếu có). \\
\hline
\textbf{Luồng phụ} & - 6.1. Sai nhiều lần (2 lần): Hệ thống hiển thị thêm mục "Gợi ý" để giúp sinh viên định hướng. \newline - 6.2. Thất bại liên tiếp (3 lần trở lên): Hệ thống hiển thị đề xuất sinh viên nên quay lại ôn tập bài học lý thuyết tương ứng. \newline - 6.3. Giải thành công: Hệ thống đánh dấu bài tập đã hoàn thành và cập nhật điểm số/tiến độ. \\
\hline
\textbf{Ngoại lệ} & - 5.a. Lỗi thực thi mã (Timeout/Crash): Hệ thống thông báo lỗi runtime và yêu cầu kiểm tra lại tối ưu mã nguồn. \newline - 5.b. Lỗi kết nối máy chủ chấm bài: Thông báo "Nộp bài thất bại, vui lòng thử lại". \\
\hline
\end{tabularx}
}

\subsubsection{Đặc tả chức năng Theo dõi tiến độ}


{
\renewcommand{\arraystretch}{1.5}
\begin{tabularx}{\textwidth}{|p{3.5cm}|X|}
\caption{Bảng đặc tả chức năng Theo dõi tiến độ}
\label{tab:bang-dac-ta-chuc-nang-theo-doi-tien-do}
\\
\hline
\textbf{Mã Use Case} & UC05 \\
\hline
\textbf{Tên Use Case} & Theo dõi tiến độ \\
\hline
\textbf{Tác nhân chính} & Sinh viên \\
\hline
\textbf{Mô tả} & Sinh viên xem tổng quan quá trình học tập của mình thông qua bảng điều khiển trung tâm (Dashboard), bao gồm các chỉ số thống kê và biểu đồ xu hướng. \\
\hline
\textbf{Tiền điều kiện} & Sinh viên đã đăng nhập thành công. \\
\hline
\textbf{Hậu điều kiện} & Sinh viên nắm rõ tình hình học tập và các mảng kiến thức cần cải thiện. \\
\hline
\textbf{Luồng chính} & 1. Sinh viên truy cập vào trang "Bảng điều khiển" (Dashboard). \newline 2. Hệ thống truy vấn dữ liệu từ CSDL và hiển thị các khối thông tin: \newline - Thống kê tổng số bài đạt, điểm số từng bài. \newline - Thống kê về điểm mạnh/điểm yếu theo thời gian. \newline - Thống kê số ngày đã làm bài (Streak). \newline 3. Hệ thống hiển thị biểu đồ xu hướng học tập (Chart) thể hiện hoạt động theo thời gian. \newline 4. Hệ thống hiển thị danh sách các bài học/bài tập gần đây nhất mà sinh viên đã tương tác. \\
\hline
\textbf{Luồng phụ} & - 2.a. Người dùng mới (Chưa có dữ liệu): Hệ thống hiển thị các chỉ số bằng 0 và khuyến khích bắt đầu bài học đầu tiên. \\
\hline
\textbf{Ngoại lệ} & - 1.a. Lỗi đồng bộ dữ liệu: Hệ thống thông báo dữ liệu đang được cập nhật và hiển thị giá trị gần nhất được lưu trong bộ nhớ tạm. \\
\hline
\end{tabularx}
}

\subsubsection{Đặc tả chức năng Chat với AI Tutor}


{
\renewcommand{\arraystretch}{1.5}
\begin{tabularx}{\textwidth}{|p{3.5cm}|X|}
\caption{Bảng đặc tả chức năng Chat với AI Tutor (Sinh viên)}
\label{tab:bang-dac-ta-chuc-nang-chat-voi-ai-tutor-sinh-vien}
\\
\hline
\textbf{Mã Use Case} & UC04 \\
\hline
\textbf{Tên Use Case} & Chat với AI Tutor \\
\hline
\textbf{Tác nhân chính} & Sinh viên \\
\hline
\textbf{Mô tả} & Sinh viên thảo luận, hỏi đáp với trợ lý AI để giải quyết các khó khăn trong quá trình học. AI Tutor cá nhân hóa phản hồi dựa trên trình độ và lịch sử lỗi của sinh viên. \\
\hline
\textbf{Tiền điều kiện} & Sinh viên đã đăng nhập. Hệ thống đã thu thập được một số dữ liệu về trình độ (Level) và lỗi thường gặp của sinh viên. \\
\hline
\textbf{Hậu điều kiện} & Cuộc hội thoại được lưu lại và các mảng kiến thức yếu của sinh viên được cập nhật vào hồ sơ phân tích. \\
\hline
\textbf{Luồng chính} & 1. Sinh viên mở cửa sổ chat hoặc chọn một cuộc trò chuyện cũ. \newline 2. Sinh viên nhập câu hỏi hoặc yêu cầu hỗ trợ về một đoạn mã nguồn. \newline 3. Hệ thống (AI Tutor) thực hiện: \newline 3.1. Phân tích hồ sơ sinh viên: Trình độ (Sơ cấp/Trung bình/Cao cấp), các ngôn ngữ đang học, các loại lỗi hay mắc phải. \newline 3.2. Áp dụng chiến lược giảng dạy (Socratic): Đưa ra các câu hỏi gợi mở, dùng phép tương đồng, dẫn dắt từng bước nhỏ thay vì đưa ra đáp án trực tiếp. \newline 3.3. Tạo tiêu đề tự động cho cuộc trò chuyện (nếu là tin nhắn đầu tiên). \newline 4. Hệ thống phản hồi dưới dạng dòng chảy dữ liệu (streaming) để sinh viên theo dõi câu trả lời thời gian thực. \newline 5. Sinh viên tiếp tục đặt câu hỏi dựa trên phản hồi của AI cho đến khi hiểu vấn đề. \newline 6. Hệ thống lưu lại toàn bộ lịch sử tin nhắn và cập nhật các thống kê về từ khóa lỗi vào CSDL phân tích. \\
\hline
\textbf{Luồng phụ} & - 3.1.a. SV Sơ cấp: AI dùng ngôn ngữ đơn giản, ví dụ gần gũi, khích lệ nhiều. \newline - 3.1.b. SV Cao cấp: AI phản hồi súc tích, yêu cầu SV phân tích các trường hợp biên và tối ưu hóa. \\
\hline
\textbf{Ngoại lệ} & - 4.a. Lỗi kết nối dịch vụ AI: Hệ thống thông báo "AI Tutor đang bận, vui lòng thử lại sau ít phút". \newline - 4.b. Tin nhắn quá dài: Hệ thống yêu cầu sinh viên rút ngắn câu hỏi để đảm bảo chất lượng phản hồi. \\
\hline
\end{tabularx}
}

\subsubsection{Đặc tả chức năng Cập nhật Knowledge log}


{
\renewcommand{\arraystretch}{1.5}
\begin{tabularx}{\textwidth}{|p{3.5cm}|X|}
\caption{Bảng đặc tả chức năng Cập nhật Knowledge log}
\label{tab:bang-dac-ta-chuc-nang-cap-nhat-knowledge-log}
\\
\hline
\textbf{Mã Use Case} & UC07 \\
\hline
\textbf{Tên Use Case} & Cập nhật Knowledge log \\
\hline
\textbf{Tác nhân chính} & Hệ thống \\
\hline
\textbf{Mô tả} & Tự động ghi lại và phân tích mọi hoạt động học tập của sinh viên để cập nhật hồ sơ năng lực và điều chỉnh trình độ học tập tương ứng. \\
\hline
\textbf{Tiền điều kiện} & Sinh viên thực hiện các tương tác (Nộp bài tập, Chat với AI) có kết quả hợp lệ hoặc lỗi. \\
\hline
\textbf{Hậu điều kiện} & Hồ sơ năng lực (Analytics) của sinh viên được cập nhật dữ liệu mới nhất. \\
\hline
\textbf{Luồng chính} & 1. Hệ thống tiếp nhận sự kiện sau mỗi lượt tương tác: \newline - Trường hợp Nộp bài Đạt: \newline \+ Hệ thống ghi nhận chủ đề bài học vừa hoàn thành vào danh sách "Kiến thức đã nắm vững". \newline \+ Đánh dấu trạng thái "Hoàn thành" cho bài tập đó trong tiến trình học tập. \newline - Trường hợp Nộp bài Lỗi/Sai: \newline \+ Hệ thống phân tích loại lỗi (Cú pháp, Logic, Runtime) và tăng bộ đếm tần suất cho loại lỗi đó trong hồ sơ. \newline 2. Hệ thống kiểm tra điều kiện nâng cấp trình độ (Level Up): \newline - Tính toán tỷ lệ vượt qua bài tập (Pass Rate) dựa trên ít nhất 5 bài nộp gần nhất. \newline - Nếu tỷ lệ đạt > 80\%, hệ thống tự động cập nhật "Ước tính trình độ" mới (ví dụ từ Sơ cấp lên Trung cấp). \newline 3. Hệ thống đồng bộ hóa các từ khóa quan trọng từ cuộc hội thoại với AI vào hồ sơ phân tích mảng kiến thức yếu/mạnh. \\
\hline
\textbf{Luồng phụ} & - 2.a. Chưa đủ số lượng bài nộp: Hệ thống tiếp tục tích lũy dữ liệu mà không thực hiện thay đổi trình độ. \\
\hline
\textbf{Ngoại lệ} & - 1.a. Lỗi cập nhật CSDL Analytics: Hệ thống thử lại sau hoặc ghi vào nhật ký lỗi chờ xử lý hậu kỳ. \\
\hline
\end{tabularx}
}


\end{document}